\documentclass[UTF8,a4paper,fontset=none]{ctexart}

\usepackage[top=2.54cm,bottom=2.54cm,left=3.17cm,right=3.17cm]{geometry}
\usepackage{fontspec}
\usepackage{xeCJK}
\usepackage{graphicx}
% Force graphics to be included even if the document/Overleaf is in draft mode
\setkeys{Gin}{draft=false}
\usepackage{booktabs}
\usepackage{array}
\usepackage{tabularx} % flexible-width tables
\usepackage{longtable}
\usepackage{etoolbox} % patch environments (tables, etc.)
\usepackage{xparse}   % redefine environments with optional args
\usepackage{amsmath,amssymb}
\usepackage{caption}
\usepackage{subcaption}
\usepackage{tocloft}
\usepackage{setspace}

% ---------------- Abstract defaults ----------------
% 用宏包统一调整摘要标题与正文为``小四号”
\usepackage{abstract}
\renewcommand{\abstractnamefont}{\xiaosi\bfseries}
\renewcommand{\abstracttextfont}{\xiaosi}

% ---------------- Keywords (摘要后的关键词) ----------------
% 用法：\Keywords{关键词1；关键词2；...}
% 要求：小四号，``关键词”三字加粗
\NewDocumentCommand{\Keywords}{m}{%
  \par\noindent{\xiaosi\textbf{关键词}：}#1\par%
}

% Better line breaking/spacing to reduce Overfull \hbox warnings
\usepackage{microtype} % subtle kerning/protrusion where supported
\usepackage{xurl}      % allow URL line breaks at more characters

\usepackage{float}
\usepackage{multirow}
\usepackage{enumitem}
\usepackage{listings}
\usepackage{inconsolata}

% hyperref is best loaded as late as possible
\usepackage{hyperref}

\lstset{
  basicstyle=\small\ttfamily,
  breaklines=true,
  breakatwhitespace=false,
  frame=single,
  numbers=left,
  numberstyle=\tiny,
  xleftmargin=2em,
  tabsize=2,
  language=MATLAB,
}

% Fonts: use portable TeX Live fonts available on Overleaf
\setmainfont{TeX Gyre Termes}
\setsansfont{TeX Gyre Heros}
\setmonofont{Inconsolata}

% CJK fonts: use Fandol (bundled with TeX Live) for portability
\setCJKmainfont{FandolSong-Regular}
\setCJKsansfont{FandolHei-Regular}
\setCJKmonofont{FandolHei-Regular}
\setCJKfamilyfont{zhsong}{FandolSong-Regular}
\setCJKfamilyfont{zhhei}{FandolHei-Regular}
\setCJKfamilyfont{zhkai}{FandolKai-Regular}
\setCJKfamilyfont{zhfs}{FandolFang-Regular}
\newcommand{\songti}{\CJKfamily{zhsong}}
\newcommand{\heiti}{\CJKfamily{zhhei}}
\newcommand{\kaishu}{\CJKfamily{zhkai}}
\newcommand{\fangsong}{\CJKfamily{zhfs}}

\hypersetup{colorlinks=false,pdfborder={0 0 0}}

\newcommand{\titlefont}{\songti}
\newcommand{\fangsongfont}{\fangsong}
\newcommand{\heifont}{\heiti}
\newcommand{\kaitifont}{\kaishu}

\newcommand{\yihao}{\fontsize{26pt}{45pt}\selectfont}
\newcommand{\erhao}{\fontsize{22pt}{45pt}\selectfont}
\newcommand{\sanhao}{\fontsize{16pt}{24pt}\selectfont}
\newcommand{\xiaosanhao}{\fontsize{15pt}{24pt}\selectfont}
\newcommand{\sihao}{\fontsize{14pt}{24pt}\selectfont}
\newcommand{\xiaosi}{\fontsize{12pt}{24pt}\selectfont}
\newcommand{\wuhao}{\fontsize{10.5pt}{18pt}\selectfont}
\newcommand{\xiaoer}{\fontsize{18pt}{18pt}\selectfont}
% ---------------- Table defaults ----------------
% 1) Font size: 五号（\wuhao）
% 2) Alignment: table 居中；tabularx 的 X 列默认居中
% 3) Width: 想``通栏到页面两端”，建议统一使用 tabularx/tabular* 并显式给 \textwidth
\AtBeginEnvironment{table}{\wuhao\centering}
\AtBeginEnvironment{tabular}{\wuhao}
\AtBeginEnvironment{tabularx}{\wuhao}
\AtBeginEnvironment{longtable}{\wuhao}

% Center X columns in tabularx (recommended for full-width tables)
\renewcommand{\tabularxcolumn}[1]{>{\centering\arraybackslash}m{#1}}
% Column helpers (so p{..} / m{..} can be centered consistently)
\newcolumntype{P}[1]{>{\centering\arraybackslash}p{#1}}
\newcolumntype{M}[1]{>{\centering\arraybackslash}m{#1}}

% Note on full-width tables:
% LaTeX's plain tabular cannot automatically stretch to \textwidth unless you
% either (a) use a stretchable column (e.g., X in tabularx) or (b) use tabular*
% with an explicit width. So for ``通栏”表格，建议统一写成：
%   \begin{tabularx}{\textwidth}{...X...}
% 或：
%   \begin{tabular*}{\textwidth}{@{\extracolsep{\fill}}...}

% Make longtable flush with both margins
\setlength{\LTleft}{0pt}
\setlength{\LTright}{0pt}


\setlength{\parindent}{2em}
\setlength{\parskip}{0pt}
% Give TeX a bit more flexibility before it overflows the margin.
\setlength{\emergencystretch}{3em}
\setlength{\tabcolsep}{4pt}
\renewcommand{\arraystretch}{1.0}
\setlength{\extrarowheight}{1pt}
\linespread{1.0}
% Keep disabled in final output.
\overfullrule=0pt
\pagestyle{plain}
\AtBeginDocument{\xiaosi}
\newenvironment{tabletext}{\begingroup\singlespacing\songti\xiaosi}{\endgroup}

% Semantic note block (use instead of manual ''\noindent 注：...'')
\NewDocumentEnvironment{Note}{+b}{%
  \par\addvspace{2pt}%
  \noindent\textbf{注：}#1\par\addvspace{2pt}%
}{}

\ctexset{
  % 一级标题：改为``1 标题”样式（数字 + 空格 + 标题），并统一小四号
  % 说明：number 使用 \thesection 等``语义编号”，以兼容 \appendix 的 A/B/… 编号。
  section = {
    name = {,},
    number = \thesection,
    aftername = {\hspace{0.5em}},
    format = \centering\heifont\xiaosi,
    beforeskip = 10pt,
    afterskip = 4pt
  },
  % 其余各级标题：统一小四号
  subsection = {
    number = \thesubsection,
    aftername = {\hspace{0.5em}},
    format = \kaitifont\xiaosi,
    beforeskip = 8pt,
    afterskip = 3pt
  },
  subsubsection = {
    number = \thesubsubsection,
    aftername = {\hspace{0.5em}},
    format = \songti\xiaosi\bfseries,
    beforeskip = 6pt,
    afterskip = 2pt
  },
  paragraph = {
    number = \theparagraph,
    aftername = {\hspace{0.5em}},
    format = \songti\xiaosi,
    beforeskip = 0pt,
    afterskip = 0pt,
    runin = false
  }
}

\DeclareCaptionFont{captionfont}{\songti\xiaosi}
\captionsetup{justification=centering,labelsep=quad,font=captionfont}
\captionsetup[table]{position=top}
\captionsetup[figure]{position=bottom}
\renewcommand{\tablename}{表}
\renewcommand{\figurename}{图}
\renewcommand{\contentsname}{{\heifont\sihao 目录}}
\renewcommand{\refname}{{\heifont\sihao 参考文献}}
\setcounter{tocdepth}{3}
\numberwithin{equation}{section}

\newcommand{\fullfig}[3]{
  \begin{figure}[H]
    \centering
    \includegraphics[width=#1\textwidth]{#2}
    \caption{#3}
  \end{figure}
}

\NewDocumentEnvironment{papertable}{o m +b}
{
  \begin{table}[H]
    \centering
    \begin{tabletext}
#3
    \end{tabletext}
    \caption{#2}
    \IfValueT{#1}{\label{#1}}
  \end{table}
}
{}

\NewDocumentEnvironment{paperlongtable}{m m +b}
{
  \begin{tabletext}
  \begin{longtable}{#1}
#3
  \caption{#2}
  \end{longtable}
  \end{tabletext}
}
{}

\begin{document}

% ---------------- Title (第一页标题) ----------------
\title{\titlefont\xiaoer 结构先验驱动的支路车流量反演与最少采样设计}
\author{}
\date{}
\maketitle
\thispagestyle{plain}

\pagestyle{plain}
\pagenumbering{arabic}
\setcounter{page}{1}

% ---------------- 摘要与关键词（第一页） ----------------
\begin{abstract}
   城市道路监测中，主路检测器仅能记录汇入总流量而无法直接观测各支路流量——这是典型的``由聚合观测反演潜在分支''的逆问题。求解过程采用以结构先验驱动的统一参数辨识框架：将支路流量的分段形态、时滞关系、信号灯开关等机理信息显式建模为可辨识约束，在约束最小二乘或稳健估计下实现支路反演；随后通过秩条件与敏感性分析将监测布设转化为面向结构参数的观测设计，在保证可重构性的前提下压缩采样规模。\par

   \vspace{1em}

   \begin{figure}[!htbp]
     \centering
     \includegraphics[width=0.5\linewidth]{Paper/paper_assets/figures/peter-robbins-cA7soLyA7i8-unsplash.jpg}
     \caption*{Photo by Peter Robbins on Unsplash}
   \end{figure}

   \textbf{问题一}针对两支路加和场景下分解不唯一的问题，将``支路1线性增长、支路2先增后减''的结构先验参数化为 $x_1(t)=pt+q$，以最小能量准则 $\min\sum(x_1^2+x_2^2)$ 附加非负约束求解。结构约束成功将无限多解转变为稳定可解的约束优化问题，分解精度达到机器精度量级（残差低于 $10^{-10}$），$R^2\approx 1$ 表明流量曲线已被精确重构。\par
   \textbf{问题二}在四支路叠加、含时滞与周期项的复杂场景中，先对主路序列一阶差分检测同向跳变间距锁定周期 $T=28$，剥离周期后再对剩余分段函数联合估计。周期项被有效剥除，非周期成分亦达到机器精度量级的重构精度。\par
   \textbf{问题三}针对信号灯对支路3的开关调制，将红绿灯建模为已知门控函数 $g(t)\in\{0,1\}$：绿灯区间内估计连续流量参数 $G_k$，红灯区间施加零流量硬约束 $\hat{x}_3(t)=0$。进一步引入可优化归零终点 $p_5$ 与下降起始点 $q_2$ 后，模型 $R^2$ 从0.7112提升至0.8699，$t\in[43,46]$ 时段最大残差降幅超过60\%。\par
   \textbf{问题四}引入加性观测误差 $\varepsilon(t)$，以粒子群优化搜索离散断点、内层约束最小二乘估计连续参数。在噪声背景下仍获得 $\mathrm{RMSE}=3.48$、$\mathrm{MAE}=2.89$、$R^2=0.9709$ 的拟合效果，且残差通过Lilliefors正态性检验（$p>0.05$）——这说明以``结构恢复''而非``逐点拟合''为目标，是含噪逆问题中保持主要形态不失真的可行路径。\par
   \textbf{问题五}将视角从反演切换至观测设计：建立``参数自由度计数→设计矩阵秩检验→关键区间覆盖验证''的三重判据。问题二最少需31个采样时刻覆盖周期全部相位（$T=28$）及分段斜率差分点；问题三改进版需22个采样时刻覆盖6个绿灯区间、支路1六阶段（含新增衰减段与零流量段）及支路2三阶段，两类方案的设计矩阵均通过满秩检验。实际部署时建议增加2~4个冗余点以提升抗噪鲁棒性。
\end{abstract}

\Keywords{结构先验；车流量反演；最少采样设计；信号灯控制；监测布设}

\clearpage

\section{问题重述}
\subsection{问题背景}
在城市道路监测中，主路监测点通常只能记录汇入后的总车流量，而难以直接观测各支路的真实流量变化。若能依据有限监测记录反推出各支路的流量函数，不仅可以为交叉口运行分析、信号优化和监测布设提供依据，也能体现统计建模中``由总体观测恢复潜在结构”的核心思想。与一般回归预测不同，本题更关注支路函数形式、关键参数及其可辨识性，因此具有明显的逆问题与结构估计特征\cite{refOD1,refOD2}。

\subsection{问题提出}
题目围绕主路单点监测条件下的支路流量分解与监测设计，设置了五个递进子问题：问题一研究最简单的两支路加和场景；问题二加入时滞、周期项和四支路叠加；问题三进一步考虑信号灯控制；问题四在问题三基础上引入观测误差；问题五则研究最少采样时刻设计。各题虽然形式不同，但都可归结为``在已知结构先验下，从主路观测反推支路函数”的统一建模任务。

\section{问题分析}
\subsection{问题一的分析}
问题一只有一个观测方程却对应两个未知支路函数——每个 $t$ 时刻的 $y(t)=x_1(t)+x_2(t)$ 构成欠定方程组，自由度为 1，解空间为无穷维。能打破这种不确定性的唯一途径是引入额外的结构约束：利用题目提供的``支路1线性增长、支路2先增后减''形态先验，将无限维的函数搜索压缩到少数几个参数构成的有限维空间，使欠定问题获得唯一解。

\subsection{问题二的分析}
问题二中周期项与各支路的分段趋势叠加在一起。若同时估计所有参数，周期基函数与分段函数的常数项之间存在共线性，可辨识性大打折扣。合理做法是利用周期项在一阶差分中呈现等间距同向跳变的独特规律，先独立锁定周期长度 $T$，再从剥离周期后的残差中估计各支路的分段参数——本质上是``特征模式识别→残差结构估计''的两步解耦策略\cite{refDyn1,refBreak}。

\subsection{问题三的分析}
问题三的特殊难点来自信号灯引入的离散控制：绿灯开启时支路 3 有流量，红灯关闭时立刻归零。若试图用一个连续的平滑分段函数同时拟合开关效应和背景流量的渐变，优化器会在开关切换点附近产生寄生振荡——因为它无法区分``流量本身的变化''和``通断控制引起的跳变''。因此必须把``何时通、何时断''作为显式的离散门控变量引入模型，使控制逻辑与连续流量的估计分属不同的决策空间\cite{refSignal1}。

\subsection{问题四的分析}
问题四主动引入了观测误差，前提假设发生变化：不再有精确穿点的必要，任何模型都无法使含噪观测的拟合残差归零。建模目标因此从``精确重构''调整为``在噪声中辨识信号的主要结构''。评价标准也从单点误差转向 RMSE、MAE、决定系数和残差分布等统计量——此时验证的是结构参数在噪声下的稳定性，而非单个样本的复现精度\cite{refRobust1}。

\subsection{问题五的分析}
问题五的视角与前四题截然不同：此前一直假定采样时刻已给定、从中反演支路参数；问题五反过来追问——如果要使反演成立，最少需要哪些时刻采样。本质上是问``哪些时间位置承载了支路参数的可辨识信息''。解决思路分为三步：遍历参数空间的自由度以确定所需信息量、构造设计矩阵并检验其列满秩性、将关键相位和拐点区间作为覆盖条件补充考量。逆问题至此转化为观测设计问题\cite{refSensor1,refSensor3,refSensor6}。

\section{模型假设}
\begin{enumerate}[itemsep=0pt,topsep=2pt,parsep=0pt,partopsep=0pt]
  \item 题目给出的分段结构先验真实有效。
  \item 主路流量可视为各支路流量在监测点处的线性叠加。
  \item 在无特别说明时，车辆从支路进入主路后的时滞按照题设固定，不考虑额外随机延迟。
  \item 历史规律可由低维分段函数与周期项刻画。
  \item 问题四中的观测偏差体现在加性误差项中。
\end{enumerate}

\section{符号说明}

\begin{papertable}{主要符号说明}
  \small
  \renewcommand{\arraystretch}{1.2}
  \begin{tabularx}{\textwidth}{>{\centering\arraybackslash}p{0.12\textwidth} >{\centering\arraybackslash}X >{\centering\arraybackslash}p{0.12\textwidth}}
    \toprule
    符号 & 含义 & 单位 \\
    \midrule
    $t$ & 离散记录时刻，$t=0,1,\dots,59$，每步 2 分钟 & 2 min \\
    $y(t)$ & 主路监测点在时刻 $t$ 的车流量观测值 & 辆 \\
    $y_{\mathrm{obs}}(t)$ & 含观测误差的主路车流量记录值 & 辆 \\
    $x_i(t)$ & 第 $i$ 条支路在时刻 $t$ 的真实车流量 & 辆 \\
    $s(t)$ & 周期性附加项或周期扰动项 & 辆 \\
    $\varepsilon(t)$ & 观测误差项 & 辆 \\
    $\tau_i$ & 第 $i$ 条支路对应的固定时滞 & 2 min \\
    $m$ & 汇入主路的支路条数 & 条 \\
    $T$ & 周期项的周期长度 & 2 min \\
    $p_i$ & 分段函数中的关键断点位置参数 & 2 min \\
    $x_i(t)$ & 支路流量，各题脚本中记作 $B_i(t)$、$f_i(t)$、$b_i(t)$、$Q_i(t)$ & 辆 \\
    \bottomrule
  \end{tabularx}
\end{papertable}

\section{模型准备}
\subsection{统一观测方程}
五个问题虽然场景不同，但都可以写成``主路总量由若干支路流量叠加形成”的统一形式。若不考虑观测误差，则有
\begin{equation}
y(t)=\sum_{i=1}^{m}x_i(t-\tau_i)+s(t),
\end{equation}
其中 $m$ 为支路条数，$\tau_i$ 为第 $i$ 条支路对应的固定时滞，$s(t)$ 表示周期性附加项；在问题四中进一步加入误差项 $\varepsilon(t)$，即可得到含噪观测模型。

\subsection{支路函数的结构表达}
\subsubsection{分段线性函数}
对具有``先增后稳”``先增后减”``多阶段变化”等特征的支路，统一采用分段线性函数进行刻画。这样既能保持参数数量可控，又便于将斜率、平台值和转折点直接解释为交通运行特征。

\subsubsection{开关控制函数}
对受红绿灯控制的支路，引入开关函数描述绿灯导通、红灯关闭的流量贡献。这样可以把连续流量变化和离散控制结构分离建模，从而避免把信号切换误判为背景流量本身的波动。

\subsubsection{周期函数}
对问题二中的周期扰动项，利用固定周期和零均值条件进行描述，并通过差分跳变位置识别其周期长度。该做法可将周期项从非周期支路中剥离出来，降低整体辨识难度。

\subsection{统一求解策略}
\subsubsection{连续参数估计}
对于斜率、截距、平台值等连续参数，统一采用约束最小二乘进行求解，并附加非负约束及必要的连续性约束，以保证流量函数具有可解释的物理意义。

\subsubsection{离散结构搜索}
对于断点位置、周期长度、首个绿灯起始时刻等离散结构参数，根据问题复杂度分别采用网格搜索或粒子群优化进行处理，再在给定结构下回代最小二乘估计连续参数。

\subsubsection{评价指标}
模型评价以 RMSE、MAE、决定系数、残差时序、设计矩阵秩和关键区间覆盖情况为主。对无噪理想情形，重点检验结构重构能力；对含噪情形，则更关注模型对主趋势和关键结构的恢复效果。

\section{问题一的模型建立及求解}
\subsection{模型建立}
\subsubsection{观测关系}
问题一中，主路流量可表示为两条支路流量之和，即 $y(t)=x_1(t)+x_2(t)$。根据题设，支路 1 线性增长，支路 2 先增后减。令支路 1 取线性形式 $x_1(t)=pt+q$，再以 $x_2(t)=y(t)-x_1(t)$ 构造支路 2，并在非负约束下最小化两支路流量平方和。

支路函数关系完整写为：
\begin{align}
  x_1(t) &= pt + q, \label{eq:q1_b1} \\
  x_2(t) &= y(t) - x_1(t) = \begin{cases}
    (7-q)+(1.5-p)t, & t\le 30, \\
    (67-q)-(0.5+p)t, & t>30.
  \end{cases} \label{eq:q1_b2}
\end{align}
其中主路 $y(t)$ 在 $t=30$（8:00）处达到峰值 52.0，前段斜率 1.5，后段斜率 $-0.5$。

\subsection{模型求解}
\subsubsection{求解准则}
问题一的关键在于：仅由主路数据并不能唯一分解出两条支路，因此必须引入结构先验。由题设可知，支路 1 的函数形式最简单，适合作为主导参数化对象；支路 2 则由主路总量减去支路 1 自动得到。若仅满足加和关系，理论上会有无限多组分解解，因此进一步引入''最小能量''准则，令
\begin{equation}
\min_{p,q}\sum_{t=0}^{59}\left[x_1^2(t)+x_2^2(t)\right],
\end{equation}
并附加 $x_1(t)\ge 0$、$x_2(t)\ge 0$ 的可行域约束，且需满足 $pt+q\le y(t)$（确保支路 2 非负）。这样得到的解具有最小波动性，能避免不必要的剧烈振荡。

求解结果得到完整的支路函数表达式：
\[
x_1(t)=0.3540t+7.0000,\qquad
x_2(t)=\begin{cases}
1.1460t, & t\le 30,\\
-0.8540t+60.0000, & t>30.
\end{cases}
\]
其中 $r_1=1.1460$、$s_1=0.0000$ 为支路 2 前段斜率与截距（$t\le 30$），$r_2=-0.8540$、$s_2=60.0000$ 为后段斜率与截距（$t>30$）。拟合残差低于 $10^{-10}$ 量级，$R^2\approx 1$，说明分解严格满足加和关系。

\subsubsection{结果图表}
\fullfig{0.92}{paper_assets/figures/q1_decomposition.png}{问题一主路与支路分解结果}
\fullfig{0.92}{paper_assets/figures/q1_branch_profiles.png}{问题一两条支路的函数形态}

\begin{papertable}[tab:q1]{问题一关键参数与统计量}
  \begin{tabularx}{\textwidth}{*{6}{>{\centering\arraybackslash}X}}
    \toprule
    参数 & 数值 & 参数 & 数值 & 指标 & 数值 \\
    \midrule
    $p$ & 0.3540 & $q$ & 7.0000 & $\overline{x_1}$ & 17.4443 \\
    $r_1$ & 1.1460 & $s_1$ & 0.0000 & $\overline{x_2}$ & 19.3057 \\
    $r_2$ & -0.8540 & $s_2$ & 60.0000 & $\max x_2$ & 34.3786 \\
    \cmidrule{4-6}
    RMSE & $\approx 10^{-12}$ & $R^2$ & $\approx 1$ & & \\
    \bottomrule
  \end{tabularx}
\end{papertable}
\begin{Note}
$r_1,s_1$ 为支路 2 前段（$t\le 30$）线性系数：$x_2(t)=r_1 t+s_1$；$r_2,s_2$ 为后段（$t>30$）系数：$x_2(t)=r_2 t+s_2$。
\end{Note}

\subsection{结果展示}
支路 1 呈平稳的线性增长趋势，表达式为 $x_1(t)=0.3540t+7.0000$，在7:00至8:58间从7.0增至27.9。支路 2 在 8:00 左右达到峰值 $34.38$ 后逐步回落至后段尾值为 $0.8540\times 59-60 \approx -9.6$（实际非负约束下取零），符合题设的先增后减特征。两支路均值分别为 17.44 和 19.31，表明两者对主路的贡献比例接近。

\subsection{结果分析与验证}
问题一说明，在低维结构假设下，仅用主路数据即可恢复支路的主要变化趋势，但恢复结果依赖于额外的辨识准则。注意到支路2在8:30后的下降段斜率为$-0.8540$，其绝对值约为主路在相同时间段斜率$(-0.5)$的1.7倍，意味着支路2的回落速率快于主路下降速率，这是由支路1的持续增长部分抵消所致。该处的''最小能量''准则实际上是 $L_2$ 正则化在无噪场景下的等价形式——在问题四的含噪情形中，同样准则将自然退化为岭回归，起到控制方差的作用。

\section{问题二的模型建立及求解}
\subsection{模型建立}
\subsubsection{时滞叠加关系}
问题二主路满足
\begin{equation}
y(t)=x_1(t-1)+x_2(t-1)+x_3(t)+s(t),
\end{equation}
其中 $x_1$ 为稳定常数，$x_2$ 为''增长-稳定-再增长''分段线性函数，$x_3$ 为''增长-稳定''函数，$s(t)$ 为周期函数。针对周期项，利用一阶差分识别跳变点，再通过同向跳变间距确定周期。

各支路与周期项的函数关系完整写为：
\begin{align}
  x_1(t) &= c_1, \label{eq:q2_f1} \\
  x_2(t) &= \begin{cases}
    k_{2e}(t+1), & -1\le t\le 24,\\
    25k_{2e}, & 24<t\le 37,\\
    25k_{2e}+k_{2l}(t-37), & 37<t\le 59,
  \end{cases} \label{eq:q2_f2} \\
  x_3(t) &= \begin{cases}
    k_3 t, & 0\le t\le t_3,\\
    k_3 t_3, & t_3<t\le 59,
  \end{cases} \label{eq:q2_f3} \\
  s(t) &= s\!\bigl(\mathrm{mod}(t,T)\bigr),\quad \sum_{i=0}^{T-1}s(i)=0. \label{eq:q2_f4}
\end{align}
其中 $t_3$ 为支路 3 的未知转折点，$T$ 为周期函数的未知周期。

\subsection{模型求解}
\subsubsection{周期识别}
问题二的可辨识性主要来自两个方面。第一，题目已明确给出三条非周期支路的形态类型，从而大幅压缩了参数自由度；第二，周期项满足零均值与固定周期的条件，可以与非周期项分离。先对主路序列作一阶差分，检测显著跳变位置；再利用同向跳变之间``隔一跳”的间距识别周期长度，最终得到 $T=28$。在周期长度固定后，剩余参数可通过约束最小二乘统一求解。

\subsubsection{参数估计}
这一处理与 Okutani 和 Stephanedes (1984) 在卡尔曼滤波框架中提出的''降维时变恢复''思路在逻辑上相通\cite{refDyn1}——两者的共同策略都是先用低维结构锚定状态主方向，再在此基底上恢复时变分量。区别在于他们的方法依赖线性状态空间模型与递推更新，而本题为一维确定性数据，可直接用闭式最小二乘完成。同时，问题二与问题三中大量出现未知转折点，可归入 Muggeo (2003) 提出的断点回归框架\cite{refBreak}，其核心是用分段线性函数拟合连续数据并在断点处施加连续性条件，与支路 2、支路 3 转折点的处理方式一致。

\subsection{结果展示}
\subsubsection{结果图表}
\fullfig{0.94}{paper_assets/figures/q2_overview.png}{问题二拟合结果、非周期分解、周期项与残差}
\fullfig{0.90}{paper_assets/figures/q2_branch_profiles.png}{问题二四个组成部分的函数对比}

\begin{papertable}[tab:q2fit]{问题二参数估计与拟合精度}
  \begin{tabularx}{\textwidth}{*{8}{>{\centering\arraybackslash}X}}
    \toprule
    $T$ & $c_1$ & $k_{2e}$ & $k_{2l}$ & $k_3$ & $t_3$ & RMSE & $R^2$ \\
    \midrule
    28 & 38.3571 & 0.6000 & 0.4000 & 1.0000 & 17 & $\approx 10^{-12}$ & $\approx 1$ \\
    \bottomrule
  \end{tabularx}
\end{papertable}

周期长度识别为 $T=28$（即56分钟），支路1为稳定常数 $c_1=38.3571$。支路2在早期（$\tau\leq 24$）斜率为 $k_{2e}=0.6$，后期（$\tau>37$）斜率 $k_{2l}=0.4$，中间段保持平台值15.0。支路3转折点 $t_3=17$（即7:34），增长斜率为1.0，之后稳定在17.0。周期项围绕零均值波动，幅值约14。拟合残差低于 $10^{-10}$ 量级，$R^2\approx 1$，模型在无噪条件下实现了参数可辨识性验证。

\subsection{结果分析与验证}
问题二是全文中结构最完整、可辨识性最强的一题：主路重构与观测值的残差仅受机器精度限制。需要重申的是，这一精度来自题目给定的严格分段结构、确定性时滞和零噪声条件，对实际交通数据而言不可复现——但它验证了生成机制完全落入模型假设空间时，参数可辨识性可以成立。利用跳变位置识别周期的方法在周期项图中得到了直观印证，这为问题五的最优采样设计提供了直接依据。

\subsubsection{7:30 与 8:30 各支路车流量数值}
\begin{papertable}[tab:q2val]{问题二 7:30 和 8:30 各支路车流量数值}
  \begin{tabularx}{\textwidth}{*{6}{>{\centering\arraybackslash}X}}
    \toprule
    时刻 & 支路1 $f_1$ & 支路2 $f_2$ & 支路3 $f_3$ & 支路4 $f_4$ & 合计 \\
    \midrule
    7:30 & 38.3571 & 9.6000 & 15.0000 & $-8.8571$ & 54.1000 \\
    8:30 & 38.3571 & 18.2000 & 17.0000 & $-8.8571$ & 64.7000 \\
    \bottomrule
  \end{tabularx}
\end{papertable}
需要注意，表 \ref{tab:q2val} 中的``支路"列给出的是各支路在对应\textbf{时刻}的流量，而模型 $y(t)=x_1(t-1)+x_2(t-1)+x_3(t)+s(t)$ 中 $x_1,x_2$ 具有 2 分钟时滞。因此在 7:30（$t=15$）时，支路1的实际流量为 $x_1(15)=c_1=38.3571$，支路2的实际流量为 $x_2(15)=0.6\times 16=9.6$（增长段），支路3的实际流量为 $x_3(15)=1\times 15=15$，周期项 $s(15)=-8.8571$，合计 54.1。按\textbf{模型公式}检验：$y(15)=x_1(14)+x_2(14)+x_3(15)+s(15)=38.3571+9.0+15.0-8.8571=53.5$，与观测值 $y(15)=53.5$ 完全吻合（残差为机器精度量级）。8:30（$t=45$）同理，按模型公式：$y(45)=x_1(44)+x_2(44)+x_3(45)+s(45)=38.3571+17.8+17.0-8.8571=64.3$，与观测值 $y(45)=64.3$ 完全吻合，验证了模型结构在无噪条件下的精确重构能力。

\section{问题三的模型建立及求解}
\subsection{模型建立}
\subsubsection{信号控制机制}
问题三的核心难点在于支路 3 受红绿灯调制，主路满足
\begin{equation}
y(t)=x_1(t-1)+x_2(t-1)+x_3(t).
\end{equation}
其中 $x_1$ 为六阶段分段函数，$x_2$ 为三阶段分段函数（下降起始点 $q_2$ 可优化），$x_3$ 在绿灯时段为正、红灯时段为零。求解采用''断点外层搜索 + 线性参数内层约束最小二乘''的思路。

各支路的函数关系完整写为：
\begin{align}
  x_1(\tau) &= \begin{cases}
    0, & 0\le\tau<p_1,\\
    a_1(\tau-p_1), & p_1\le\tau<p_2,\\
    a_1(p_2-p_1)-\beta_1(\tau-p_2), & p_2\le\tau<p_3,\\
    S_1, & p_3\le\tau<p_4,\\
    \displaystyle S_1\cdot\frac{p_5-\tau}{p_5-p_4}, & p_4\le\tau<p_5,\\
    0, & \tau\ge p_5,
  \end{cases} \label{eq:q3_b1} \\
  x_2(\tau) &= \begin{cases}
    a_2\tau+b_{20}, & 0\le\tau\le35,\\
    a_2\cdot35+b_{20}=S_2, & 35<\tau<q_2,\\
    S_2-d_2(\tau-q_2), & q_2\le\tau\le59,
  \end{cases} \label{eq:q3_b2} \\
  x_3(t) &= \begin{cases}
    G_k, & t\in\text{第 }k\text{ 个绿灯区间},\\
    0, & t\in\text{红灯区间},
  \end{cases} \label{eq:q3_b3}
\end{align}
其中 $\tau=t-1$ 为时滞标度，$p_1,p_2,p_3,p_4,p_5$ 为支路 1 的五个断点（$p_5$ 为归零终点），$a_1,\beta_1,S_1$ 为支路 1 的线性参数（由连续性约束 $a_1(p_2-p_1)-\beta_1(p_3-p_2)=S_1$ 关联），$a_2,b_{20},S_2,d_2$ 为支路 2 的线性参数，$q_2$ 为支路 2 的下降起始点（原固定 $\tau=47$，现作为可优化参数），$G_1,\dots,G_7$ 为 7 个绿灯区间的常值流量参数。信号灯周期为 9 个时间单位（绿 5 红 4），首个绿灯起始于 $t=3$。

\subsection{模型求解}
\subsubsection{求解难点}
问题三相比问题二出现了两个新的统计难点。其一，支路 3 的贡献不是简单平滑变化，而是被红绿灯强行切换为''有流量/无流量''的交替模式，这会导致主路序列产生更强的非平稳特征。其二，支路 1 的多阶段模型（六阶段）包含五个未知断点，若断点选择稍有偏差，就会把由信号灯引起的跳变误判为某条支路本身的变化。因此本题拟合效果较问题二明显下降是合理现象，不应被简单视为模型失效。为改善模型在 $t=43\sim 46$ 时段的系统性高估，在原始五阶段基础上做了三项改进：(i) 将支路 1 归零终点从固定 $t=59$ 改为可优化参数 $p_5$，使衰减斜率更灵活；(ii) 将支路 2 下降起始点从固定 $\tau=47$ 改为可优化参数 $q_2$，允许提前进入下降段；(iii) 去掉了 $d_2\ge 0.01$ 的硬约束，改为 $d_2\ge 0$。

\subsubsection{求解依据}
Webster (1958) 关于信号配时的经典研究\cite{refSignal1}系统给出了单点交叉口最优周期、绿信比与延误的关系，但其分析建立在固定相位配时、均匀到达的假设之上。本题与之不同之处在于：到达过程非均匀且未知（由三条支路叠加而成），检测器也仅有主路总量而无分相位计数。这意味着信号切换与流量变化同时耦合在观测序列中，分离二者的难度高于经典信号配时设计场景。因此，本题重点不在于追求极低残差，而在于能否恢复绿灯开启时段内的主要流量贡献与背景流量形态。

\subsubsection{结果图表}
\fullfig{0.94}{paper_assets/figures/q3_overview.png}{问题三主路拟合、支路分解、信号灯控制与残差}
\fullfig{0.90}{paper_assets/figures/q3_branch_profiles.png}{问题三各支路函数及关键断点}

\begin{papertable}{问题三断点参数与拟合指标（改进后）}
  \begin{tabularx}{\textwidth}{*{9}{>{\centering\arraybackslash}X}}
    \toprule
    $p_1$ & $p_2$ & $p_3$ & $p_4$ & $p_5$ & $q_2$ & $a_1$ & $S_1$ & $R^2$ \\
    \midrule
    4 & 14 & 28 & 34 & 37 & 51 & 6.0591 & 44.4930 & 0.8699 \\
    \bottomrule
  \end{tabularx}
\end{papertable}

参数含义如下：支路1（$b_1$）为六阶段函数——$[0,p_1)$零流量、$[p_1,p_2)$以斜率$a_1=6.0591$增长（峰值为60.59）、$[p_2,p_3)$以斜率$\beta_1=1.1498$线性下降至$S_1=44.49$、$[p_3,p_4)$保持稳定值$S_1$、$[p_4,p_5)$以斜率$\gamma_1=14.83$快速衰减至零、$[p_5,59]$保持零流量。支路2（$b_2$）为三阶段函数——$[0,35]$增长（斜率$a_2=1.5643$）、$(35,q_2)$稳定（$S_2=54.75$）、$[q_2,59]$以斜率$d_2=3.3968$下降。支路3（$b_3$）在7个绿灯区间内取值分别为21.91、7.84、14.91、12.09、13.65、19.95、12.46，红灯时段为零。拟合$R^2=0.8699$说明改进后模型能够解释主路约87\%的方差波动，相比原始模型（$R^2=0.7112$）提升了约22个百分点。

\subsection{结果展示}
问题三的难度高于问题二，因为信号灯使得观测值呈现明显的开关调制效应。支路 3 在绿灯区间的脉冲式贡献被有效建模为分段常值，其拟合结果基本反映了信号控制对流量分配的影响。支路1在早高峰时段（$t \in [0,30]$）承担了主要的流量输送任务，支路2则在上午时段维持相对稳定的背景流量，支路3的开关式贡献反映了交叉口信号配时的周期性切换。

\subsection{结果分析与验证}
改进后的模型在 $t=43\sim 46$ 时段的最大残差从 $-28.0$ 降至 $-10.7$，$R^2$ 从 0.7112 提升至 0.8699。关键在于三项改动使模型获得了更大的灵活性：$p_5=37$ 使支路 1 在 $[34,37)$ 以斜率 14.83 快速归零（原模型从 $t=43$ 缓慢衰减至 $t=59$，斜率仅约 3.0），$q_2=51$ 使支路 2 的稳定期延长至 $\tau=51$ 才开始下降（原模型固定 $\tau=47$ 即开始下降，导致 $S_2$ 被压低），去掉了 $d_2$ 下界后下降斜率 $d_2=3.40$ 完全由数据驱动。当前模型仍存在以下局限。第一，支路 3 在绿灯期内仍为分段常值假设，每个绿灯区间内流量被压缩为一个常数 $G_k$，无法反映绿灯期间流量爬升或回落的动态变化。第二，在 $t=48\sim 52$ 时段（绿灯6），观测值从 44 跃升至 82 后又快速回落至 38（$t=53$），这种锯齿状波动难以被分段线性结构完全捕获。第三，断点搜索以离散步长（1个时间单位 = 2 分钟）进行，$p_1\sim p_5$ 和 $q_2$ 的时间精度仅到 2 分钟量级。总体而言，改进后的模型在处理离散控制与连续流量耦合场景时取得了明显优于原始模型的拟合效果。

\subsubsection{7:30 与 8:30 各支路车流量数值}
\begin{papertable}[tab:q3val]{问题三 7:30 和 8:30 各支路车流量数值}
  \begin{tabularx}{\textwidth}{*{5}{>{\centering\arraybackslash}X}}
    \toprule
    时刻 & 支路1 $b_1$ & 支路2 $b_2$ & 支路3 $b_3$ & 合计 \\
    \midrule
    7:30 & 59.4408 & 23.4642 & 7.8446 & 90.7496 \\
    8:30 & 0.0000 & 54.7499 & 0.0000 & 54.7499 \\
    \bottomrule
  \end{tabularx}
\end{papertable}
7:30时，信号灯处于第二绿灯区间 $G_2(t \in [12,16])$，支路3贡献 $G_2=7.84$，支路1处于下降段末端（$p_2=14 < t-1=14 = p_2$，$b_1=60.5906$），支路2处于增长段（$t-1=14 \leq 35$，$b_2=1.5643\times 14+0.0=21.90$），合计约90.75与观测值92.79接近。8:30时，$b_1$ 因 $p_5=37 < t-1=44$ 已归零，$b_2$ 处于稳定段（$35 < t-1=44 < q_2=51$）贡献 54.75，$b_3$ 处于红灯区间为零，合计 54.75 高于观测值 44.0。残差 -10.75 较改进前的 -28.3 已有显著缩小，但仍高于观测值约 20\%。此偏差的剩余部分来自支路 3 在 8:30 附近红灯周期内仅被赋予零流量而未能捕获残余交通流的过渡效应，同时支路 2 的稳定段 $S_2=54.75$ 在这一时段偏高，是模型在高流量区间与低流量区间之间的最优折中选择。

\section{问题四的模型建立及求解}
\subsection{模型建立}
\subsubsection{含噪观测方程}
问题四在问题三的基础上加入设备误差项，
\begin{equation}
y_{\mathrm{obs}}(t)=x_1(t)+x_2(t)+x_3(t)+\varepsilon(t).
\end{equation}
各支路函数关系为：
\begin{align}
  x_1(t) &= \begin{cases}
    0, & t<\theta_1,\\
    k_1(t-\theta_1), & \theta_1\le t<\theta_2,\\
    C_1=k_1(\theta_2-\theta_1), & \theta_2\le t<\theta_3,\\
    C_1-k_2(t-\theta_3), & \theta_3\le t<\theta_4,\\
    0, & t\ge\theta_4,
  \end{cases} \label{eq:q4_q1} \\
  x_2(t) &= \begin{cases}
    \alpha t+\beta, & 0\le t\le 17,\\
    \alpha\cdot17+\beta=C_2, & 17<t\le 35,\\
    C_2-\gamma(t-35), & 35<t\le59,
  \end{cases} \label{eq:q4_q2} \\
  x_3(t) &= \begin{cases}
    C_{3,k}, & t\in\text{第 }k\text{ 个绿灯区间},\\
    0, & \text{红灯区间},
  \end{cases} \label{eq:q4_q3}
\end{align}
其中 $\theta_1,\theta_2,\theta_3,\theta_4$ 为支路1梯形函数的四个断点（隐含连续性约束 $C_1=k_1(\theta_2-\theta_1)=k_2(\theta_4-\theta_3)$），$\theta_g$ 为首个绿灯起始相位（表 \ref{tab:q4} 中记为 $\tau_g$）；$\alpha,\beta,\gamma,C_2$ 为支路2参数；$C_{3,k}$ 为第 $k$ 个绿灯区间内的常值流量。采用符号 $\theta$ 代替 $\tau$ 以避免与符号表中时滞 $\tau_i$ 冲突。考虑到断点与首个绿灯时刻属于离散参数，而各支路斜率与平台值属于连续参数，采用''外层粒子群优化搜索离散结构，内层最小二乘估计连续参数''的双层策略，从而兼顾搜索能力与计算效率。

\subsection{模型求解}
\subsubsection{误差口径}
问题四不再追求对观测值的严格穿点，而是把未解释部分全部纳入 $\varepsilon(t)$。因此本题的评价重点也从``是否精确重构”转为``是否在合理误差范围内恢复主要结构”。以均方根误差、平均绝对误差和决定系数作为主指标，同时观察残差随时间的分布、符号与收敛曲线，以判断误差是否主要来源于设备偏差而非模型失真。

\subsubsection{算法依据}
这一设定与交通计数中存在不一致时的处理方法在目标上有相似之处。例如，Nie 等 (2004) 在处理带容量约束的交通分配问题时，通过引入惩罚项和松弛变量来吸收观测异常值\cite{refRobust1}。但区别在于他们的方法面向网络级 OD 矩阵估计，需要借助多源计数冗余来识别异常位置；而本题仅有一条主路序列，没有冗余信息可供交叉验证，因此采用 PSO$+$最小二乘的双层策略在结构层面吸收误差，而非在残差层面做鲁棒回归。

\subsubsection{结果图表}
\fullfig{0.94}{paper_assets/figures/q4_overview.png}{问题四拟合结果、鲁棒分解、残差与收敛曲线}
\fullfig{0.90}{paper_assets/figures/q4_branch_profiles.png}{问题四三条支路的恢复流量曲线}

\begin{papertable}[tab:q4]{问题四关键参数与误差指标}
  \begin{tabularx}{\textwidth}{*{8}{>{\centering\arraybackslash}X}}
    \toprule
    $\theta_1$ & $\theta_2$ & $\theta_3$ & $\theta_4$ & $\theta_g$ & RMSE & MAE & $R^2$ \\
    \midrule
    3.8392 & 28.4494 & 48.3973 & 57.6438 & 1.3822 & 3.48 & 2.89 & 0.9709 \\
    \bottomrule
  \end{tabularx}
\end{papertable}

\subsection{结果展示}
含噪条件下，模型对主路整体趋势的跟踪能力保持良好，PSO 收敛过程较稳定。与问题二和问题三相比，问题四更贴近真实监测场景——观测误差无法回避，而双层优化框架能够在误差背景下恢复主要的支路结构信息。

\subsection{结果分析与验证}
模型的 $R^2=0.9709$ 说明方法能捕获主路约97\%的方差变化，这是对''结构恢复能力''的直接度量，即支路流量的主要波形和时域定位未被噪声破坏。PSO 收敛曲线约在150代后达到稳定。RMSE=3.48 约为观测值量级（17～88）的 6\%，说明误差主要来源于设备引入的加性噪声而非模型结构失真——如果误差来自模型形式错误（如漏掉某条支路），残差会呈现与信号相关的系统偏差，其幅值往往与信号本身同量级。该场景下的 RMSE 远小于信号量级，指向的是零均值随机误差。Lilliefors 检验接受残差正态性假设（$p>0.05$），进一步支持了''残差为设备噪声''的解释，也意味着过拟合风险较低——因为过拟合通常导致残差模式与信号结构相关，而非近似独立同分布。

\subsubsection{7:30 与 8:30 各支路车流量数值}
\begin{papertable}[tab:q4val]{问题四 7:30 和 8:30 各支路车流量数值}
  \begin{tabularx}{\textwidth}{*{7}{>{\centering\arraybackslash}X}}
    \toprule
    时刻 & 支路1 $Q_1$ & 支路2 $Q_2$ & 支路3 $Q_3$ & 合计 & 观测值 & 残差 \\
    \midrule
    7:30 & 17.49 & 18.18 & 20.31 & 55.99 & 55.99 & 0.00 \\
    8:30 & 38.58 & 0.00 & 0.00 & 38.58 & 35.42 & -3.16 \\
    \bottomrule
  \end{tabularx}
\end{papertable}
\begin{Note}
由于 PSO 算法存在随机性，每次运行所得具体数值略有不同，但总体结构和误差水平保持稳定。7:30 时，支路 1 处于增长段（$\theta_1=3.84 < 15 < \theta_2=28.45$），支路 2 处于第一段线性增长（$t=15<17$），且信号灯处于绿灯相位（相位偏移约 4.62，处于绿灯区间），因此三条支路均有贡献，合计约 56 辆。8:30 时，支路 1 处于稳定段（$\theta_2=28.45<45<\theta_3=48.40$），支路 2 已过减少段后降至零，信号灯处于红灯相位，因此仅支路 1 贡献约 38.58 辆。残差 $-3.16$ 的绝对值小于 RMSE$=3.48$，说明该偏差处于模型平均误差的可接受范围内。
\end{Note}

\section{问题五的模型建立及求解}
\subsection{模型建立}
\subsubsection{设计目标}
问题五研究在已知支路函数结构的前提下，监测设备至少应在哪些时刻采样才能恢复完整函数。以参数自由度、设计矩阵秩与典型采样覆盖关系为核心进行分析，并分别给出问题二和问题三的关键采样集合。

\subsection{模型求解}
\subsubsection{可辨识性条件}
问题五的本质是可辨识性分析。若某一模型包含 $k$ 个独立线性参数，则采样所对应的设计矩阵至少应达到秩 $k$，否则参数无法唯一确定。对问题二而言，周期项 $f_4$ 引入了 $T$ 个相位值（受零均值约束后为 $T-1$ 个独立参数），加上支路1的 $c_1$、支路2的 $k_{2e},k_{2l}$、支路3的 $k_3$，共计 $T+3$ 个线性参数（$T=28$ 时即31个），因此采样重点在于覆盖全部周期相位和分段线性区间的关键转折点。对问题三的改进版而言，支路1（2个线性参数+5个断点$p_1\sim p_5$）、支路2（3个线性参数+1个可优化下降起始点$q_2$）、支路3（7个绿灯区间参数）共12个线性参数和6个离散结构参数，采样重点在于覆盖每个绿灯区间、各分段区间以及新增的衰减段与零流量段。两类问题的采样逻辑并不相同。

\subsubsection{采样设计依据}
问题五的采样设计本质上属于交通监测中的传感器选址与部分可观测性分析范畴。Yang等\cite{refSensor1}从OD矩阵估计的可靠性出发讨论了最少检测器配置条件，其秩条件思路与设计矩阵满秩要求有直接对应关系；Chow和Lo\cite{refSensor3}从信息论角度选择最优检测器位置，最大信息量原则与覆盖关键区间以保证参数可辨识性的思想一致；Ng\cite{refSensor6}从节点层面推导了无需路径枚举的传感器布设条件，其图论结构与分段区间覆盖约束逻辑相通。借鉴上述工作，这里采用"参数自由度 + 设计矩阵秩 + 关键区间覆盖"三重判据来确定最少采样集：首先依据参数计数确定采样数下界，再构建设计矩阵并检验是否满秩，最后验证所选时刻是否覆盖了所有关键相位、转折点和控制区间。

\subsubsection{结果图表}
\fullfig{0.94}{paper_assets/figures/q5_sampling_design.png}{问题五中问题二与问题三的最少采样时刻设计}

\begin{papertable}{问题五参数自由度与采样规模}
  \small
  \setlength{\tabcolsep}{3pt}
  \begin{tabularx}{\textwidth}{*{6}{>{\centering\arraybackslash}X}}
    \toprule
    模型 & 线性参数数 & 离散参数数 & 设计矩阵秩 & 最少采样数 & 备注 \\
    \midrule
    问题二 & 31 & 1 & 31 & 31 & 周期项主导 \\
    问题三 & 12 & 6 & 12 & 22 & 含 b1 衰减段与零流量段 \\
    \bottomrule
  \end{tabularx}
\end{papertable}

具体而言，问题二的最小采样集包含以下31个时刻：为覆盖周期 $T=28$ 的全部相位，需在 $t=0,1,\ldots,26$ 各取1个点（共27个时刻覆盖相位 $0$ 至 $26$，相位 $27$ 由零均值约束间接确定）；另外为识别支路2的斜率参数 $k_{2e}$ 和 $k_{2l}$，需在 $t=34$（差分消去周期后确定 $k_{2e}$）和 $t=55$（确定 $k_{2l}$）额外采样；加之覆盖支路1常数项和支路3转折点所需的2个冗余采样。问题三（改进版）的最小采样集包含以下22个时刻：每个绿灯区间中部各取1点（$t=4,12,17,26,31,36,49$），支路2增长段2点（$t=5,20$）、稳定段1点（$t=40$）、减少段1点（$t=55$），支路1六阶段覆盖点（增长段$t=7,15$、减少段$t=28$、稳定段$t=38$、衰减段$t=36$、零流量段$t=52$），以及边界覆盖点（$t=2,3,35,47,59$）。相比原始五阶段模型（21点），新增 $t=36$ 覆盖新增的衰减段 $[p_4,p_5)$。该采样集的设计矩阵经检验满秩，可唯一确定全部线性参数。

\subsection{结果展示}
问题二采样点均匀分布在 $0$ 至 $26$ 的相位上，覆盖了周期 $T=28$ 的全部相位需求，同时额外补充了两个用于识别分段斜率的差分采样点。问题三的采样更强调绿灯时段与分段结构的联合覆盖：每个绿灯区间内均有一个采样点以确保对应流量参数的可辨识性；支路1的六个阶段区间（含新增的衰减段和零流量段）和支路2的三个阶段区间也分别被采样点覆盖，其中 $t=36$ 专门针对衰减-零流量边界。两类方案的设计矩阵均通过了满秩检验。

\subsection{结果分析与验证}
考虑到真实监测数据必然含噪，实际部署时不宜严格按照理论最少采样点执行。较稳妥的做法是在理论最少采样集基础上增加 2 至 4 个冗余点，优先布置在周期跳变附近、信号灯切换附近以及分段斜率变化附近。以问题二为例，建议在 $t=14$（周期下跳点）、$t=42$（周期下跳点）和 $t=26$（周期平稳段）各增加一个冗余采样；以问题三为例，建议在支路1的各转折点 $p_1=4$、$p_2=14$、$p_3=28$、$p_4=34$、$p_5=37$ 附近各增加一个采样点以增强转折点识别的鲁棒性。这样既能提升数值稳定性，也更符合统计建模中''设计留有冗余''的经验原则。

\section{模型的评价与推广}
\subsection{模型的评价}
综合五题结果：第一，已知支路变化结构的条件下，主路单点监测数据足以恢复支路流量函数——问题一（两支路分解，$R^2\approx 1$）和问题二（四支路+周期项，机器精度量级）均验证了这一点。第二，当交通控制机制或观测误差引入额外不确定性后，模型仍然可以恢复主要结构信息，但拟合精度明显下降——问题三改进后 $R^2=0.8699$ 和问题四的 $R^2=0.9709$ 分别反映了信号灯调制和观测噪声对重构精度的影响。通过引入可优化的 $p_5$ 和 $q_2$，问题三的 $R^2$ 较原始模型提升了约22个百分点。第三，最优采样设计的核心是根据模型自由度和控制结构选择关键时刻——问题二需要31个时刻覆盖周期相位，问题三（改进版）需要22个时刻覆盖六阶段支路和绿灯区间。

\subsection{模型的推广}
"结构先验 + 约束优化"框架存在若干推广方向，但每项推广都会引入当前模型所不曾涉及的困难，需要相应改造。

第一，更多支路或更复杂拓扑结构（如多级汇入或环形交叉口）的车流量分解。当前框架依赖于"主路观测 = 各支路流量直接叠加"的假设，多级汇入场景下这一假设将失效——上游汇入的车流可能在下游交叉口再次分流。此时需要将叠加关系替换为路径-路段关联矩阵，结构先验也需要从支路层面提升到路径层面来定义，最小采样设计的秩判据则需相应修正为关联矩阵的可逆性条件。环形交叉口的情形更为棘手，因为分流比例本身可能随时间变化，原本固定的分段参数会变成时变量。

第二，更为精细的控制机制建模。引入自适应信号配时或多相位控制后，信号灯切换不再是一个固定的周期表，而是依赖于实时排队长度。此时"信号灯相位作为固定已知输入"这一前提不再成立，切换时刻本身成为待估计参数。这意味着在问题三的框架中，绿灯区间参数不再由外部给定，而需要与流量参数联合求解——这会显著增加参数空间的非凸性，优化算法容易陷入局部极值。

第三，传感器布设策略优化。参数自由度+设计矩阵秩+关键区间覆盖的三重判据适用于任意分段线性或周期结构，但若支路函数包含非线性成分（如指数增长段），设计矩阵将不再是常数矩阵而依赖于参数本身。此时秩条件无法在采样前独立检验，需要采用自适应采样策略——先粗采获取参数估计，再根据估计结果决定补充采样时刻。这一流程将原本的静态设计问题转化为序贯决策问题。

在应用层面，当道路检测设备发生故障或缺失部分时段数据时，上述框架可从稀疏采样反推完整流量函数，这有助于降低城市道路交通监测系统的建设与运维成本。

以下三个方向在理论上值得进一步探索，但各自面临不同的障碍。

（1）推广至非线性函数形式。当前的分段线性模型若替换为分段多项式或样条，参数自由度将大幅增加，且多项式的基函数之间容易产生共线性，设计矩阵的条件数会急剧恶化。此时不仅需要增加采样点数量以对抗方差膨胀，还需引入Tikhonov正则化或Lasso来抑制过拟合——这等价于在当前秩判据之外又附加了一个正则化路径选择的维度。

（2）引入贝叶斯估计框架。当前框架中，分段断点$p_i$和信号灯切换时刻等结构参数以确定性约束（等式或不等式）形式纳入优化。若改为概率先验，需要解决的不仅是"把优化换成采样"，更在于先验分布的合理选择：交通流量的分段断点通常集中在典型时段（如早高峰起始），这一先验信息应编码为稀疏性先验（如Laplace先验或Horseshoe先验），而不同参数的先验依赖结构——如$p_i$与$q_j$在实际中往往正相关——也会使后验采样的计算成本急剧上升。

（3）发展在线参数更新算法。当前所有求解均为批处理方式——将全部观测时刻的数据一次性代入优化。若改为递推估计（如递推最小二乘或卡尔曼滤波），需要解决两个新问题：一是分段断点的在线检测与更新——断点在递推中无法像批处理那样"回头看"全局数据来定位，而需要依赖累计误差的变点检测机制；二是当新增数据触发分段结构变化时，如何在不重启整个优化的情况下将历史信息平顺地转移到新结构下。这两个问题在现有的递推估计算法框架下仍缺乏通用解法。

\clearpage
\begingroup
\singlespacing
\xiaosi
\raggedright
\begin{thebibliography}{99}
\setlength{\itemsep}{0pt}
\setlength{\parsep}{0pt}
\bibitem{ref2} Treiber M, Kesting A. Traffic Flow Dynamics: Data, Models and Simulation[M]. Berlin: Springer, 2013.
\bibitem{refOD1} van Zuylen H J, Willumsen L G. The most likely trip matrix estimated from traffic counts[J]. Transportation Research Part B, 1980, 14(3): 281--293.
\bibitem{refOD2} Cascetta E, Nguyen S. A unified framework for estimating or updating origin/destination matrices from traffic counts[J]. Transportation Research Part B, 1988, 22(6): 437--455.
\bibitem{refDyn1} Okutani I, Stephanedes Y J. Dynamic prediction of traffic volume through Kalman filtering theory[J]. Transportation Research Part B, 1984, 18(1): 1--11.
\bibitem{refBreak} Muggeo V M R. Estimating regression models with unknown break-points[J]. Statistics in Medicine, 2003, 22(19): 3055--3071.
\bibitem{refRobust1} Nie Y, Zhang H M, Lee D-H. Models and algorithms for the traffic assignment problem with link capacity constraints[J]. Transportation Research Part B, 2004, 38(4): 285--312.
\bibitem{refSignal1} Webster F V. Traffic signal settings[R]. Road Research Technical Paper No. 39, London: HMSO, 1958.
\bibitem{refSensor1} Yang H, Iida Y, Sasaki T. An analysis of the reliability of an origin-destination trip matrix estimated from traffic counts[J]. Transportation Research Part B, 1991, 25(5): 351--363.
\bibitem{refSensor3} Chow A H F, Lo H P. An information-theoretic sensor location model for traffic origin-destination demand estimation applications[J]. Transportation Science, 2010, 44(2): 254--273.
\bibitem{refSensor6} Ng M. Synergistic sensor location for link flow inference without path enumeration: A node-based approach[J]. Transportation Research Part B, 2012, 46(6): 781--788.
\end{thebibliography}
\endgroup

\clearpage
\section*{附录}
\addcontentsline{toc}{section}{附录}
\appendix
\section{支撑材料目录}
支撑材料包括以下文件，存放于 \path{Solutions/} 与 \path{Paper/paper_assets/} 目录中：

\begin{enumerate}[itemsep=0pt,topsep=2pt,parsep=0pt,partopsep=0pt]
  \item 问题一：\path{Solutions/Q1.m}（约束最小二乘优化脚本 ㊣）、\path{Paper/paper_assets/tables/q1_summary.csv}（参数与结果汇总）、\path{Paper/paper_assets/figures/q1_decomposition.png}、\path{Paper/paper_assets/figures/q1_branch_profiles.png}（结果图）
  \item 问题二：\path{Solutions/Q2.m}（周期检测+约束最小二乘脚本，含辅助函数 \texttt{fit\_}）、\path{Paper/paper_assets/tables/q2_summary.csv}（参数与支路流量汇总）、\path{Paper/paper_assets/figures/q2_overview.png}、\path{Paper/paper_assets/figures/q2_branch_profiles.png}（结果图）
  \item 问题三：\path{Solutions/Q3.m}（网格搜索+约束最小二乘脚本，含 $p_5$ 和 $q_2$ 改进）、\path{Paper/paper_assets/tables/q3_summary.csv}（参数与支路流量汇总）、\path{Paper/paper_assets/figures/q3_overview.png}、\path{Paper/paper_assets/figures/q3_branch_profiles.png}（结果图）
  \item 问题四：\path{Solutions/Q4.m}（PSO+约束最小二乘脚本，含辅助函数 \texttt{evaluate\_}、\texttt{decode\_}、\texttt{compute\_} 等）、\path{Paper/paper_assets/tables/q4_summary.csv}（参数与支路流量汇总）、\path{Paper/paper_assets/figures/q4_overview.png}、\path{Paper/paper_assets/figures/q4_branch_profiles.png}（结果图）
  \item 问题五：\path{Solutions/Q5.m}（采样集设计与矩阵秩分析脚本，含辅助函数 \texttt{build\_}、\texttt{q2\_}）、\path{Paper/paper_assets/tables/q5_summary.csv}（设计结果汇总）、\path{Paper/paper_assets/figures/q5_sampling_design.png}（结果图）
\end{enumerate}

\section{MATLAB 代码}
本文所有模型求解均在 MATLAB R2024a 环境下完成。各问题的核心代码（\path{Solutions/Q1.m} 至 \path{Solutions/Q5.m}）如下所示。

\subsection{问题一主脚本}
\lstinputlisting{../Solutions/Q1.m}

\subsection{问题二主脚本}
\lstinputlisting{../Solutions/Q2.m}

\subsection{问题三主脚本}
\lstinputlisting{../Solutions/Q3.m}

\subsection{问题四主脚本}
\lstinputlisting{../Solutions/Q4.m}

\subsection{问题五主脚本}
\lstinputlisting{../Solutions/Q5.m}

\end{document}
